\documentclass{beamer}

\usetheme{Madrid}
\usecolortheme{seahorse}
\usepackage[utf8]{inputenc}
\usepackage[spanish]{babel}
\usepackage{listings}
\usepackage{xcolor}

% --- Configuración del estilo del código C# ---
\lstdefinestyle{CSharpStyle}{
  language=[Sharp]C,
  basicstyle=\scriptsize\ttfamily,
  keywordstyle=\color{blue}\bfseries,
  stringstyle=\color{teal},
  commentstyle=\color{gray}\itshape,
  showstringspaces=false,
  breaklines=true,
  frame=single,
  rulecolor=\color{gray!30},
  tabsize=2,
  captionpos=b,
  numbers=left,
  numberstyle=\tiny\color{gray},
}

\title{Uso de List (Lista) en C\# para el Desarrollo de Videojuegos con Unity}
\author{COM200-Programacion Avanzada}
\date{\today}

\begin{document}

\begin{frame}
\titlepage
\end{frame}

%------------------------------------------------
\begin{frame}{Concepto general}
\begin{itemize}
    \item Una \textbf{List (lista)} en C\# es una estructura de datos dinámica que almacena elementos de forma secuencial.
    \item A diferencia de un arreglo (Array) tradicional, su capacidad de memoria puede \textbf{incrementarse o reducirse dinámicamente} en tiempo de ejecución.
    \item En el contexto de la programación, se conceptualiza como un contenedor secuencial donde los objetos se estructuran mediante posiciones numéricas (índices).
    \item En el motor Unity, esta estructura resulta de particular utilidad para:
    \begin{itemize}
        \item Diseñar sistemas de inventario dinámicos y escalables.
        \item Administrar conjuntos de entidades activas en la escena (e.g., enemigos, proyectiles).
        \item Almacenar nodos de navegación (waypoints) para sistemas de Inteligencia Artificial.
    \end{itemize}
\end{itemize}
\end{frame}

%------------------------------------------------
\begin{frame}[fragile]{Declaración de un List}
Para implementar una lista en C\#, se requiere la inclusión previa del espacio de nombres correspondiente a las colecciones genéricas:
\begin{lstlisting}[style=CSharpStyle]
using System.Collections.Generic;
\end{lstlisting}

Posteriormente, se procede con la declaración e inicialización de la estructura, especificando el tipo de dato subyacente que albergará:
\begin{lstlisting}[style=CSharpStyle]
List<GameObject> inventario = new List<GameObject>();
\end{lstlisting}
\end{frame}

%------------------------------------------------
\begin{frame}[fragile]{Operaciones Add() e Insert() — Inserción de elementos}
El método \textbf{Add()} adjunta un elemento al final de la secuencia de la lista.
\begin{lstlisting}[style=CSharpStyle]
inventario.Add(espada);
\end{lstlisting}

El método \textbf{Insert()} introduce un elemento en una posición (índice) específica, desplazando los elementos subsecuentes.
\begin{lstlisting}[style=CSharpStyle]
inventario.Insert(0, pocionDeVida); // Inserción en la posición inicial
\end{lstlisting}

Implementación en el contexto de Unity:
\begin{lstlisting}[style=CSharpStyle]
public void RecolectarElemento(GameObject nuevoElemento)
{
    inventario.Add(nuevoElemento);
    Debug.Log($"Elemento registrado. Total en estructura: {inventario.Count}");
}
\end{lstlisting}
\end{frame}

%------------------------------------------------
\begin{frame}[fragile]{Operaciones Remove() y RemoveAt() — Eliminación de elementos}
El método \textbf{Remove()} localiza el elemento especificado y suprime su primera aparición. Retorna un valor booleano (\texttt{true}) si la operación de eliminación es exitosa.
\begin{lstlisting}[style=CSharpStyle]
bool exito = inventario.Remove(espadaRota);
\end{lstlisting}

El método \textbf{RemoveAt()} elimina el elemento ubicado exactamente en el índice proporcionado como argumento.
\begin{lstlisting}[style=CSharpStyle]
if (inventario.Count > 0)
{
    inventario.RemoveAt(0); // Supresión del primer elemento de la secuencia
}
\end{lstlisting}
\end{frame}

%------------------------------------------------
\begin{frame}[fragile]{Acceso por Índice — Consulta y Modificación}
A diferencia de las estructuras tipo Pila o Cola, la Lista facilita el acceso aleatorio directo a los elementos mediante su índice (con base cero) utilizando la notación de corchetes \texttt{[]}.

\begin{lstlisting}[style=CSharpStyle]
// Acceso de lectura al tercer elemento (índice 2)
GameObject elementoActivo = inventario[2];
Debug.Log("Objeto referenciado: " + elementoActivo.name);

// Sobrescritura de un elemento existente en la estructura
inventario[1] = escudoMejorado;
\end{lstlisting}

\textit{Nota:} Es imperativo garantizar que el valor del índice sea estrictamente menor a la propiedad \texttt{Count} para prevenir el lanzamiento de la excepción \texttt{ArgumentOutOfRangeException}.
\end{frame}

%------------------------------------------------
\begin{frame}[fragile]{Otras operaciones relevantes}
\begin{itemize}
    \item \textbf{Clear()}: Libera la estructura eliminando la totalidad de sus elementos.
    \begin{lstlisting}[style=CSharpStyle]
    inventario.Clear();
    \end{lstlisting}

    \item \textbf{Count}: Retorna la magnitud actual de la estructura (constituye una propiedad de la clase, no un método).
    \begin{lstlisting}[style=CSharpStyle]
    int magnitudTotal = inventario.Count;
    \end{lstlisting}

    \item \textbf{Contains()}: Verifica la existencia de un elemento específico dentro de la colección.
    \begin{lstlisting}[style=CSharpStyle]
    if (inventario.Contains(llaveMaestra))
        Debug.Log("Condición lógica de validación cumplida.");
    \end{lstlisting}
\end{itemize}
\end{frame}

%------------------------------------------------
\begin{frame}[fragile]{Ejemplo práctico: Controlador de Inventario}
\begin{lstlisting}[style=CSharpStyle]
using UnityEngine;
using System.Collections.Generic;

public class ControladorInventario : MonoBehaviour
{
    public List<string> elementos = new List<string>();

    public void RegistrarElemento(string identificador)
    {
        elementos.Add(identificador);
        Debug.Log("Registro exitoso: " + identificador);
    }

    public void ProcesarElemento(int indice)
    {
        if (indice >= 0 && indice < elementos.Count)
        {
            Debug.Log("Procesando entidad: " + elementos[indice]);
            elementos.RemoveAt(indice); // Eliminación tras su procesamiento
        }
        else
        {
            Debug.LogWarning("Excepción prevenida: Índice fuera de los límites.");
        }
    }
}
\end{lstlisting}
\end{frame}

%------------------------------------------------
\begin{frame}{Ventajas de la implementación de List en Unity}
\begin{itemize}
    \item \textbf{Flexibilidad de memoria:} La capacidad de almacenamiento se adapta dinámicamente según los requerimientos de ejecución.
    \item \textbf{Acceso Aleatorio:} Posibilita la lectura o modificación de los elementos en tiempo constante ($O(1)$) mediante su identificador posicional.
    \item \textbf{Riqueza de funciones:} Proporciona métodos integrados para la búsqueda, ordenamiento (\texttt{Sort}) y manipulación eficiente de la información.
    \item \textbf{Iterabilidad:} Facilita la iteración sistemática mediante el uso de estructuras de control de flujo como \texttt{for} o \texttt{foreach}.
\end{itemize}
\end{frame}

%------------------------------------------------
\begin{frame}{Comparativa de Estructuras de Datos Lineales}
\begin{table}[]
\centering
\begin{tabular}{|l|l|l|}
\hline
\textbf{Estructura} & \textbf{Paradigma de acceso} & \textbf{Aplicación típica en simulación} \\ \hline
Stack & LIFO (Último en entrar, primero en salir) & Procesamiento inverso, historiales \\ \hline
Queue & FIFO (Primero en entrar, primero en salir) & Gestión de turnos, colas de ejecución \\ \hline
List & \textbf{Acceso directo (Aleatorio)} & \textbf{Sistemas dinámicos, repositorios} \\ \hline
\end{tabular}
\end{table}
\end{frame}

%------------------------------------------------
\begin{frame}{Conclusión}
\begin{itemize}
    \item La clase genérica \textbf{List} en C\# constituye una de las colecciones más versátiles y fundamentales en la arquitectura de software.
    \item Integra la eficiencia de acceso posicional característica de los arreglos estáticos con la adaptabilidad del redimensionamiento de memoria dinámico.
    \item En el ecosistema de Unity, el dominio de esta estructura resulta indispensable para el diseño de arquitecturas escalables, la gestión de entidades dinámicas y la administración óptima de conjuntos de datos de longitud variable durante el ciclo de vida de la aplicación.
\end{itemize}
\end{frame}

\end{document}